\begin{env}[5.1.12]
\label{IV.2.5.1.12-fr}
Sean $X$ un preesquema, $\mathcal F$ un $\mathcal O_X$-módulo
cuasicoherente \emph{de tipo finito}, $S=\operatorname{Supp}(\mathcal F)$ su
soporte, que es \emph{cerrado} en $X$ $(\mathbf{0}_{\mathrm I},5.2.2)$. Si, para todo
$x\in X$, se considera $\operatorname{Supp}(\mathcal F_x)$ como una parte
cerrada del esquema local $\operatorname{Spec}(\mathcal O_x)$, se tiene, por
definición $(\mathbf{0},16.1.7)$
$\dim(\mathcal F_x)=\dim(\operatorname{Supp}(\mathcal F_x))$; pero se tiene
\[
\operatorname{Supp}(\mathcal F_x)
=S\cap\operatorname{Spec}(\mathcal O_{X,x})
=\operatorname{Spec}(\mathcal O_{S,x})
\]
donde, en el último término, $S$ designa un subpreesquema cerrado arbitrario de
$X$ que tiene a $S$ como espacio subyacente. Si se observa que las componentes
irreducibles de $\operatorname{Spec}(\mathcal O_{S,x})$ son las
intersecciones de $\operatorname{Spec}(\mathcal O_{X,x})$ con las componentes
irreducibles de $S$ que contienen a $x$, y corresponden biunívocamente a estas
últimas, se ve que
\begin{equation}
\dim(\mathcal F_x)=\dim(\mathcal O_{S,x})
\tag{5.1.12.1}\label{IV.2.5.1.12.1-fr}
\end{equation}
en virtud de (5.1.1.1); se deduce también de (5.1.2) que se tiene
\begin{equation}
\dim(\mathcal F_x)=\operatorname{codim}(\overline{\{x\}},S)
\tag{5.1.12.2}\label{IV.2.5.1.12.2-fr}
\end{equation}
si $X$ es localmente noetheriano.

Se dice que $\mathcal F$ es \emph{equidimensional en el punto} $x\in X$ si
$\mathcal F_x$ es un $\mathcal O_{X,x}$-módulo equidimensional,
es decir, $(\mathbf{0},16.1.7)$, si $\operatorname{Supp}(\mathcal F_x)$ es
equidimensional como parte cerrada de
$\operatorname{Spec}(\mathcal O_{X,x})$; equivale a decir que el anillo
$\mathcal O_{S,x}$ es \emph{equidimensional}.

Se llama \emph{dimensión de} $\mathcal F$ y se denota $\dim(\mathcal F)$ la
dimensión del \emph{soporte} $\operatorname{Supp}(\mathcal F)$; de (5.1.4) y
(5.1.12.1) se deduce que se tiene
\begin{equation}
\dim(\mathcal F)=\sup_{x\in X}\dim(\mathcal F_x).
\tag{5.1.12.3}\label{IV.2.5.1.12.3-fr}
\end{equation}
Si $X=\operatorname{Spec}(A)$ es un esquema afín y si
$\mathcal F=\widetilde M$, donde $M$ es un $A$-módulo de tipo finito, se tiene
$\dim(\mathcal F)=\dim(M)$ en virtud de $(\mathbf{0},16.1.7)$ y (5.1.4).
\end{env}

\begin{proposition}[5.1.13]
\label{IV.2.5.1.13-fr}
Sean $X$ un preesquema, $\mathcal F$ un $\mathcal O_X$-módulo
cuasicoherente de tipo finito, $x$ un punto de $X$, $x'$ una generización de
$x$ en $X$; se tiene entonces
\begin{equation}
\dim(\mathcal F_{x'})\leq\dim(\mathcal F_x)
\tag{5.1.13.1}\label{IV.2.5.1.13.1-fr}
\end{equation}
\end{proposition}

\begin{proof}
Esto resulta inmediatamente de (5.1.12.1) y de las definiciones.
\end{proof}
